%% evenend-sample-options.tex：evenendの見本（オプションごとの組み方の違い）
%% マニュアル（evenend-ja.tex）が貼り込む．オプションを\evenendsetupで換えながら
%% 同じ材料を組む．貼り込むページの番号を evenend-sample-options.pg に書き出す
\documentclass[twocolumn,paper=a6,fontsize=8pt]{jlreq}
\usepackage{evenend}
\setlength\columnseprule{0.4pt}
\pagestyle{empty}
\newcommand\para{吾輩は猫である．名前はまだ無い．どこで生れたかとんと見当がつかぬ．何でも薄暗いじめじめした所でニャーニャー泣いていた事だけは記憶している．\par}
\newcount\parascnt
\newcommand\paras[1]{\parascnt=#1\relax\loop\ifnum\parascnt>0 \para\advance\parascnt-1 \repeat}
\makeatletter
%% 貼り込むページの番号：\samplepage{名前}を置いたページの番号を書き出す
\newwrite\evs@out
\immediate\openout\evs@out=\jobname.pg
\newcommand\samplepage[1]{\write\evs@out{\noexpand\@namedef{evs@#1}{\thepage}}}
%% floatgridの見本：本文の行の位置（行送りの格子）を薄い線で示す
\newif\ifevs@grid
\AddToHook{shipout/background}{\ifevs@grid
  \begingroup
    \unitlength=1sp
    \@tempdima=\dimexpr1in+\topmargin+\headheight+\headsep+\topskip\relax
    \@tempcnta\z@
    \@whilenum\@tempcnta<\numexpr\textheight/\baselineskip\relax\do{%
      \put(\number\dimexpr1in+\oddsidemargin\relax,
           -\number\dimexpr\@tempdima+\@tempcnta\baselineskip\relax){%
        \color[gray]{.7}\rule{\textwidth}{.2pt}}%
      \advance\@tempcnta\@ne}%
  \endgroup\fi}
\makeatother
\usepackage{xcolor}
\begin{document}

%% ---- footnote：page（全幅の脚注）とcolumn（段の下端の脚注）
\newcommand\fnsample[1]{%
  \samplepage{#1}%
  \section{脚注}
  \paras{2}
  脚注を付ける\footnote{1つ目の脚注．footnote=pageでは版面の全幅で版面の下端に置く．}．
  \paras{3}
  ここにも脚注\footnote{2つ目の脚注．}．
  \paras{3}
  最後のページの脚注\footnote{3つ目の脚注．}．
  \paras{1}
  \samplepage{#12}%
  \evenendclearpage}
\fnsample{footnote-page}
\evenendsetup{footnote=column}\setcounter{section}{0}\setcounter{footnote}{0}
\fnsample{footnote-column}
\evenendsetup{footnote=page}

%% ---- onecolumn：\onecolumnの直前のページ
\newcommand\onesample{%
  \section{段組の終わり}
  \paras{4}
  \onecolumn
  1段組の続き．
  \twocolumn}
\setcounter{section}{0}
\samplepage{onecolumn-true}\onesample
\evenendsetup{onecolumn=false}\setcounter{section}{0}
\samplepage{onecolumn-false}\onesample
\evenendsetup{onecolumn=true}

%% ---- enable：文書末のページ（ここでは\evenendclearpageで同じ組み方を示す．
%% enable=falseの見本はこの文書の最後のページ）
\newcommand\docendsample{%
  \evenendcolumns{2}%
  \section{文書の終わり}
  \paras{4}}
\setcounter{section}{0}
\samplepage{enable-true}\docendsample
\evenendclearpage

%% ---- blockskip：段数を切り換えたときの段組の間
\newcommand\skipsample{%
  \evenendcolumns{2}%
  \section{2段}
  \paras{2}
  \evenendcolumns{3}
  \paras{3}
  \evenendclearpage}
\setcounter{section}{0}
\samplepage{blockskip-2}\skipsample
\evenendsetup{blockskip=4\baselineskip}\setcounter{section}{0}
\samplepage{blockskip-4}\skipsample
\evenendsetup{blockskip=2\baselineskip}

%% ---- blockmin：切り換えた後に残った高さが小さいとき
\newcommand\minsample{%
  \evenendcolumns{2}%
  \section{2段}
  \paras{4}
  \evenendcolumns{3}
  \paras{2}
  \evenendclearpage}
\setcounter{section}{0}
\samplepage{blockmin-5}\minsample
\evenendsetup{blockmin=12\baselineskip}\setcounter{section}{0}
\samplepage{blockmin-12}\minsample
\evenendsetup{blockmin=5\baselineskip}

%% ---- headingskip：切り換えた後の段組が見出しで始まるとき
\newcommand\headsample{%
  \evenendcolumns{2}%
  \section{2段}
  \paras{2}
  \evenendcolumns{3}
  \section{3段}
  \paras{3}
  \evenendclearpage}
\setcounter{section}{0}
\samplepage{headingskip-column}\headsample
\evenendsetup{headingskip=block}\setcounter{section}{0}
\samplepage{headingskip-block}\headsample
\evenendsetup{headingskip=column}

%% ---- spanfootnote：段をまたぐ下の図と脚注（footnote=columnのとき）
\evenendsetup{footnote=column}
\newcommand\spansample{%
  \evenendcolumns{3}%  3段のうち2段にまたがる図
  \section{段をまたぐ図}
  \paras{1}
  脚注\footnote{1段目の脚注．}を付ける．
  \begin{spanfigure}[b]{2}
    \centering\fbox{\parbox[c][30pt][c]{.9\linewidth}{\centering 2段にまたがる図}}
    \caption{段をまたぐ図}
  \end{spanfigure}
  \paras{2}
  脚注\footnote{2段目の脚注．長めの脚注．長めの脚注．}を付ける．
  \paras{1}
  \evenendclearpage}
\setcounter{section}{0}\setcounter{footnote}{0}
\samplepage{spanfootnote-above}\spansample
\evenendsetup{spanfootnote=bottom}\setcounter{section}{0}\setcounter{footnote}{0}
\samplepage{spanfootnote-bottom}\spansample
\evenendsetup{spanfootnote=above,footnote=page}

%% ---- floatgrid：フロートの後の行の位置（薄い線は行送りの位置）
\makeatletter
\newcommand\gridsample{%
  \evenendcolumns{2}%
  \global\evs@gridtrue
  \section{フロートと行の位置}
  \paras{1}
  \begin{figure}[t]
    \centering\fbox{\parbox[c][25pt][c]{.8\linewidth}{\centering 図}}
    \caption{高さが行送りの倍数でない図}
  \end{figure}
  \paras{3}
  \evenendclearpage
  \global\evs@gridfalse}
\makeatother
\setcounter{section}{0}\setcounter{figure}{0}
\samplepage{floatgrid-true}\gridsample
\evenendsetup{floatgrid=false}\setcounter{section}{0}\setcounter{figure}{0}
\samplepage{floatgrid-false}\gridsample
\evenendsetup{floatgrid=true}

%% ---- enable=false：文書末のページを揃えない
\evenendsetup{enable=false}\setcounter{section}{0}
\samplepage{enable-false}\docendsample
\end{document}
